\usepackage[utf8]{inputenc}
\usepackage[T1]{fontenc}
\usepackage[brazil]{babel}
\usepackage{csquotes}
\usepackage{microtype}
\usepackage{lmodern}
\usepackage{indentfirst}
%\usepackage{parskip} % sem indentações
\usepackage[onehalfspacing]{setspace}
\usepackage[lmargin=3cm,tmargin=3cm,rmargin=2cm,bmargin=2cm,headheight=28pt]{geometry} % [showframe]
\usepackage[output-decimal-marker={,},per-mode=symbol]{siunitx}

\usepackage{lipsum}
\setlipsumdefault{1}

\usepackage{graphicx}
\usepackage{multicol,multirow}
\usepackage{booktabs}
\usepackage{titlesec}
\usepackage{caption}
\usepackage{subcaption}
\usepackage{hyperref}
\usepackage{tocloft}
\usepackage{enumitem}
\usepackage{adjustbox}
\usepackage{float}
\floatstyle{plaintop}
\restylefloat{table}

\usepackage{silence} % silencia erros e warnings como abaixo:
\WarningFilter{biblatex}{File 'brazilian-abnt-abnt.lbx' not found}

%%%%%%%%%%%%%%%%%%%%%%%%%%%%%%%%%%%%%%%%%%%%%%%%%%%%%%%%%%%%
%  Adiciona pontos nas seções no sumário (classe article)  %
%%%%%%%%%%%%%%%%%%%%%%%%%%%%%%%%%%%%%%%%%%%%%%%%%%%%%%%%%%%%

\renewcommand\cftsecfont{\normalfont}
\renewcommand\cftsecpagefont{\normalfont}
\renewcommand{\cftsecleader}{\cftdotfill{\cftdotsep}}

%%%%%%%%%%%%%%%%%%%%%%%%%%%%%%%%
%  referencias bibliográficas  %
%%%%%%%%%%%%%%%%%%%%%%%%%%%%%%%%

% Documentação biblatex-abnt
% http://linorg.usp.br/CTAN/macros/latex/contrib/biblatex-contrib/biblatex-abnt/doc/biblatex-abnt.pdf
\usepackage[backend=biber,
style=abnt,             % estilo alfabético
%style=abnt-numeric,    % estilo numérico
giveninits,             % Apenas iniciais do primeiro nome do autor
justify,                %justifica as referências ao invés de alinhar à esquerda
noslsn,                 % oculta opções como s.l s.n
backref                 % Especifica as páginas em que cada entrada foi citada
]{biblatex}

\usepackage{fancyhdr}
\fancyhf{}
\pagestyle{fancy}
\lhead{\includegraphics[height=0.8\headheight]{C:/Users/victo/repos/latex-templates/figuras/ucam-header.png}}
\fancyhead[RO]{\thepage}
\renewcommand{\headrulewidth}{0.8pt}

\newlist{enumalpha}{enumerate}{1}   % Cria um novo ambiente enumerate com letras
\setlist[enumalpha]{label=\alph*)}  % pode-se criar vários sem alterar o original

\usepackage{color}   %May be necessary if you want to color links

\hypersetup{
    colorlinks=true, %set true if you want colored links
    linktoc=all,     %set to all if you want both sections and subsections linked
    linkcolor=blue,  %choose some color if you want links to stand out
    citecolor=black,
    urlcolor=blue
}

